% =============================================================================
% Section 6 — Conclusions and Future Work
% PINN-ONB01 Phase 1 manuscript (IJHMT submission)
% Target length: ~500 words (single flowing narrative)
% =============================================================================

\section{Conclusions and Future Work}\label{sec:conclusions}

This study realized the five contributions outlined in
\cref{sec:introduction}. We curated an open dataset of \num{1361}
boiling-curve points, $82$ ONB labels, and $49$ surface descriptor cards
from seven literature sources spanning four fluids. We built a
surface-conditioned PINN with $24\,005$ parameters, combining a
nine-channel encoder, a ten-token categorical embedding, and a FiLM
modulation layer. We trained the network with a five-term composite loss
comprising the conduction residual, the natural-convection BC, an ONB
data term, a soft Hsu prior, and $R_a$/$\theta$ monotonicity
regularizers. We performed an inverse $r_c$ study over $48$ surfaces,
which revealed a Simpson-type reversal in $R_a$--$r_c$ across fluids.
We adopted a three-level verification protocol covering code accuracy,
physical trends, and deep-ensemble UQ ($K=10$).

These contributions yielded the following empirical findings. On $n=77$
held-out ONB labels, the forward RMSE was \SI{3.42}{K} with MAE
\SI{2.21}{K} and $R^2=+0.44$, a \SI{52.6}{\percent} reduction relative to
the strongest classical baseline (Basu et al., \SI{7.21}{K}) and the
only configuration with positive $R^2$. The improvement was larger on
refrigerant subsets than on the full set. RMSE reductions on R-134a and
R-123 reached \SIrange{65}{67}{\percent} (\SI{1.18}{K} and \SI{2.35}{K},
respectively), compared with \SI{46}{\percent} for water. The five
canonical monotonic trends were recovered with mean $|\rho|=0.90$, and
pressure and subcooling reached $|\rho|=1.00$. The Hsu-analytical
inverse produced mean $r_c$ \SI{3.21}{\micro\meter} with
\SI{60}{\percent} of the $48$ surfaces falling within
$[1,100]\,\si{\micro\meter}$. Stratifying by fluid revealed a sign
reversal, from $\rho=-0.27$ overall to $+0.30$ on water and $-0.77$ on
R-134a. The ensemble returned $\bar{\sigma}_{\mathrm{total}}=\SI{4.09}{K}$
(epistemic-dominated, $\bar{\sigma}_{\mathrm{epi}}=\SI{3.85}{K}$) with
$98.7\%$ empirical $\pm 2\sigma$ coverage against the nominal $95\%$.
Training took \SIrange{17}{20}{min} per replica on a single Apple~M1 CPU.

Several limitations frame the interpretation of these results. Many
surfaces have only a single ONB observation, so within-surface $q''$
dependence must be learned through the Hsu prior. Fluid coverage is
uneven, and FC-77 was excluded because its properties are not available
in CoolProp. No co-located SEM/AFM cavity measurements are available,
so the inverse $r_c$ is validated only through Hsu self-consistency.
The PDE residual on dense interior grids exceeds the $10^{-3}$ threshold
despite the training-time term converging to $\mathcal{O}(10^{-8})$,
reflecting a regression-vs-solver trade-off.

Future work will address these gaps along four directions. First, the
surface-card corpus will be enlarged with experimental partners to
include \numrange{5}{10} surfaces with direct SEM/AFM measurements and
$+50$ ONB labels from three-to-five additional sources. Second, the
framework will be extended to forced-convection subcooled boiling. The
encoder and Hsu/monotonicity regularizers will be retained, with
Reynolds number and inlet subcooling added as conditioning inputs.
Third, the 1D kernel will be promoted to two dimensions to capture
lateral spreading and instability, a prerequisite for biphilic patterns
whose ONB localizes at hydrophobic islands. Fourth, a heteroscedastic
aleatoric head will be added on top of the ensemble. The
$\sigma_{\mathrm{epi}}$ ranking will then drive an active-learning loop
with experimental collaborators.

To support these extensions and broader community use, the dataset,
trained ensemble, and training/evaluation scripts have been released
under an open-source license, providing a reproducible baseline for
future ONB studies.
